\documentclass[11pt]{article}
\usepackage[T1]{fontenc}
\usepackage{amsmath,amssymb,amsthm}
\usepackage[margin=1in]{geometry}
\usepackage{array,longtable}
\usepackage{xurl}
\usepackage[hidelinks]{hyperref}
\newtheorem{theorem}{Theorem}[section]
\newtheorem{proposition}[theorem]{Proposition}
\newtheorem{lemma}[theorem]{Lemma}
\newtheorem{corollary}[theorem]{Corollary}
\theoremstyle{definition}
\newtheorem{definition}[theorem]{Definition}
\theoremstyle{remark}
\newtheorem{remark}[theorem]{Remark}
\newcommand{\Tr}{\operatorname{Tr}}
\newcommand{\im}{\operatorname{im}}
\newcommand{\rk}{\operatorname{rank}}
\newcommand{\eps}{\varepsilon}
\newcommand{\Del}{\Delta}
\title{Unramified algebras over a discrete field\\
\large A constructive product decomposition}
\author{}
\date{31 October 2026}

\begin{document}
\maketitle

\begin{abstract}
We prove that a finitely presented, formally unramified algebra over a
discrete field is a finite product of separable monogenic algebras.
The diagonal idempotent supplies flatness transfer, regularity and
trace duality. A bounded search then gives either a global primitive
element, obtained from a tensor separator, or a finite coefficient
field over which an explicit idempotent splitting produces the required
factors. The proof uses finite linear algebra and finite searches.
\end{abstract}

\part{The mathematical argument}

\section{Statement and conventions}

All rings are commutative with identity, and algebra maps preserve
identity. The base $k$ is a nonzero \emph{discrete field}: equality in
$k$ is decidable and every nonzero element has an inverse. Algebras
over $k$ may be zero. A module is \emph{finite} if it has a finite
spanning family; a \emph{finite basis} includes unique coordinates.
We distinguish these two notions throughout the proof.
An \emph{explicitly finite} set comes with a finite enumeration and
decidable equality; this is stronger than finite generation of a module.

\begin{definition}
An $R$-algebra $C$ is \emph{formally unramified} if, for every
$R$-algebra $T$ and ideal $N$ with $N^2=0$, the map
\[
 \operatorname{Hom}_{R\text{-alg}}(C,T)
 \longrightarrow \operatorname{Hom}_{R\text{-alg}}(C,T/N)
\]
is injective. A monic polynomial $f\in k[T]$ is \emph{separable} if
$f'$ is invertible in $k[T]/(f)$, equivalently if
$uf+vf'=1$ for some $u,v\in k[T]$.
\end{definition}

\begin{theorem}\label{thm:main}
Let $C$ be a finitely presented, formally unramified $k$-algebra.
There exist finitely many monic separable polynomials $f_i\in k[T]$
and a $k$-algebra isomorphism
\[
 C\simeq\prod_i k[T]/(f_i).
\]
In particular, $C$ has a finite $k$-basis and is finite \'etale over $k$.
\end{theorem}

The factors in this statement are monogenic \'etale algebras.
Irreducibility of the $f_i$ is not required. The polynomial $1$ is
allowed: its quotient is the zero algebra, and its derivative is a
unit in that ring. Such factors can be omitted when forming the product.

We first construct the diagonal idempotent and deduce regularity.
An elementary finiteness argument then provides a basis and trace
duality. The final step has two branches, chosen by a bounded finite
calculation; it requires no decision whether $k$ is finite or infinite.

\section{The diagonal idempotent, flatness and regularity}

Write
\[
 D=C\otimes_R C,\qquad
 \mu:D\longrightarrow C,\quad \mu(a\otimes b)=ab,
 \qquad \Del a=a\otimes1-1\otimes a.
\]

\begin{lemma}\label{lem:diagonal}
If $C$ is a finitely generated, formally unramified $R$-algebra,
there is an element $\eps\in D$ such that
\[
 \mu(\eps)=1,\qquad \eps\Del a=0\quad(a\in C),
 \qquad \eps^2=\eps.
\]
Moreover $\ker\mu=(1-\eps)D$.
\end{lemma}
\begin{proof}
Put $J=\ker\mu$. If $x_1,\ldots,x_r$ generate $C$ as an algebra,
then $J$ is generated by $\Del x_1,\ldots,\Del x_r$. Indeed, after
these differences have been killed, the two maps from $C$ coincide,
and multiplication identifies the resulting quotient with $C$.

The maps $C\rightrightarrows D/J^2$ agree after passage to $D/J$.
The kernel $J/J^2$ is square-zero, so formal unramifiedness makes
the maps equal in $D/J^2$. Hence $J=J^2$.

Here is the determinant argument for this finite idempotent ideal.
Choose generators $u_1,\ldots,u_r$ of $J$ and write
$u_i=\sum_j a_{ij}u_j$ with $a_{ij}\in J$. The adjugate identity for
$I-(a_{ij})$ shows that
\[
 \eps=\det(I-(a_{ij}))
\]
annihilates every $u_i$. Also $\eps\equiv1\pmod J$.
Thus $\eps J=0$, $1-\eps\in J$, and $\eps(1-\eps)=0$.
For $u\in J$, the equality $u=(1-\eps)u$ gives
$J=(1-\eps)D$. The remaining assertions follow immediately.
\end{proof}

\begin{lemma}[Flatness transfer]\label{lem:flatness}
Suppose $R\to C$ has the diagonal element of Lemma~\ref{lem:diagonal}.
If a $C$-algebra $T$ is flat over $R$, it is flat over $C$.
\end{lemma}
\begin{proof}
The $C$-module $C\otimes_R T$ is flat by base change.
Let $\eps_T$ be the image of $\eps$ in this algebra. Multiplication
$C\otimes_R T\to T$ has the section
\[
 t\longmapsto\eps_T(1\otimes t).
\]
The equality $\eps_T(c\otimes1)=\eps_T(1\otimes c)$ makes this
section $C$-linear, and $\mu_T(\eps_T)=1$ proves the section equation.
A module retract of a flat module is flat.
\end{proof}

\begin{proposition}\label{prop:regular}
A finitely generated, formally unramified algebra $C$ over $k$ is
von Neumann regular:
\[
 \text{for every }a\in C\text{ there is }b\in C\text{ with }a=a^2b.
\]
Consequently $C$ is reduced.
\end{proposition}
\begin{proof}
Every module over a discrete field is flat. The equational criterion
proves this directly: in a finite relation $\sum_i\alpha_i t_i=0$,
either all coefficients vanish or one coefficient $\alpha_p$ is
invertible. In the second case, the vectors
$e_i-\alpha_i\alpha_p^{-1}e_p$ generate the relations needed by that
criterion. This argument requires no choice of a basis of the module.

Apply Lemma~\ref{lem:flatness} to $C/(a)$. This quotient is flat over
$C$. Tensoring the injection $(a)\hookrightarrow C$ with $C/(a)$
therefore gives an injection
\[
 (a)/(a)^2\longrightarrow C/(a).
\]
The map is zero, so $(a)=(a)^2$ and $a=a^2b$ for some $b$.
Iterating this equation gives $a=a^m b^{m-1}$ for $m\geq1$.
Thus every nilpotent element is zero.
\end{proof}

\section{Finite spanning families, bases and minimal polynomials}

\begin{definition}
Call a ring \emph{algebraically zero-dimensional} if
\[
 \forall x\;\exists N\geq0,\ b\qquad x^N(1-xb)=0.
\]
\end{definition}
Regular rings have this property with $N=1$. We use the displayed
equation itself, without any definition through prime ideals.

\begin{lemma}[Integral descent]\label{lem:descent}
If $A\hookrightarrow B$ is integral and $B$ is algebraically
zero-dimensional, then so is $A$.
\end{lemma}
\begin{proof}
For $x\in A$, choose $N,b$ with $x^N=bx^{N+1}$ in $B$, and a monic
equation
\[
 b^h+c_{h-1}b^{h-1}+\cdots+c_0=0,\qquad c_i\in A.
\]
Induction gives $x^N=b^j x^{N+j}$. Hence, for $0\leq j\leq h$,
$b^j x^{N+h}=x^{N+h-j}$. Multiplying the monic equation by
$x^{N+h}$ yields
\[
 x^N(1+c_{h-1}x+\cdots+c_0x^h)=0.
\]
This is the required equation in $A$, by injectivity. A degree-zero
monic equation gives $1=0$, when the assertion is immediate.
\end{proof}

\begin{proposition}\label{prop:finite}
A finitely generated, algebraically zero-dimensional $k$-algebra
is finite as a $k$-module. If it is finitely presented as an algebra,
it has a finite $k$-basis.
\end{proposition}
\begin{proof}
For the first assertion, induct on the number of algebra generators
$x_1,\ldots,x_r$. With no generators, $1$ spans the algebra.
For $r>0$, choose $N,b$ with $x_1^N(1-x_1b)=0$ and write
$b=P(x_1,\ldots,x_r)$. The polynomial
\[
 F=X_1^N(1-X_1P)
\]
vanishes at the generators and is nonzero: its coefficient at
$X_1^N X_2^0\cdots X_r^0$ is $1$.

Choose $M\geq2$ larger than all exponents occurring in $F$. Make the
invertible triangular substitution
\[
 X_i=Y_i+T^{M^i}\quad(1\leq i<r),\qquad X_r=T.
\]
The top $T$-degree of a monomial $X^\alpha$ is
$\alpha_r+\sum_{i<r}\alpha_iM^i$. These degrees are distinct, by
uniqueness of base-$M$ expansions. Consequently the highest
$T$-coefficient of the transformed $F$ is a nonzero scalar in $k$.
Dividing by it gives a monic equation in $T$. A nonzero constant
relation gives the zero algebra and can be handled immediately.

Let $A'$ be the image subalgebra generated by the new $r-1$ variables,
and let $t$ denote the image of $T$. The inverse substitution gives
$C=A'[t]$. If the monic equation has degree $d$, then
$1,t,\ldots,t^{d-1}$ span $C$ over $A'$. In particular $C$ is integral
over $A'$. Here one can use the determinant trick: multiplication by
an element, written on a finite spanning family containing $1$, gives
a matrix whose monic characteristic polynomial annihilates that
element. Lemma~\ref{lem:descent} makes $A'$ zero-dimensional.
Induction makes $A'$ finite over $k$, and combining the two spanning
families proves the assertion for $C$.

For the basis assertion, write
$C=k[X_1,\ldots,X_r]/(g_1,\ldots,g_s)$. The determinant trick and
the finite spanning family give monic annihilators $p_i(X_i)$ of
the original generators. The algebra
\[
 Q_0=k[X_1,\ldots,X_r]/(p_1(X_1),\ldots,p_r(X_r))
\]
has its explicit bounded-monomial basis. Its quotient by the images
of the $g_j$ is $C$. The ideal generated by these images is spanned
over $k$ by their products with the bounded basis monomials of $Q_0$.
Gaussian elimination on this finite list of vectors constructs a basis
of $C$. This step uses the original finite presentation; it does not
require a finite presentation of the intermediate image subalgebras.
\end{proof}

\begin{lemma}\label{lem:reduced-regular}
A reduced $k$-algebra which is finite as a $k$-module is regular.
\end{lemma}
\begin{proof}
For $a$ choose a monic annihilator $p(T)$ by the determinant trick.
Decidable equality in $k$ allows us to write
$p(T)=T^s(c+Th(T))$ with $c\ne0$.
If $s=0$, then $c+ah(a)=0$. If $s>0$, the element
$a(c+ah(a))$ is nilpotent and is therefore zero.
In either case $ca+a^2h(a)=0$, so
$a=a^2(-c^{-1}h(a))$.
\end{proof}

\begin{lemma}\label{lem:minimal}
Let $C$ have a finite $k$-basis, and let $x\in C$. There is a monic
polynomial $f$ of least degree with $f(x)=0$. Evaluation induces an
injection $k[T]/(f)\hookrightarrow C$, whose image is $k[x]$.
\end{lemma}
\begin{proof}
Use Gaussian elimination to test successively whether $x^m$ is in
the span of $1,x,\ldots,x^{m-1}$. A monic annihilator supplies a finite
bound for this search. The first dependence gives $f$. There is no
nonzero polynomial of smaller degree vanishing at $x$: divide such
a polynomial by its nonzero leading coefficient to obtain an earlier
dependence. Division by the monic polynomial $f$ now shows that the
kernel of evaluation is precisely $(f)$.
\end{proof}

\section{Trace duality and the monogenic separability criterion}

For a finite-basis algebra $C$, write $\Tr(a)$ for the trace of
multiplication by $a$.

\begin{proposition}\label{prop:trace}
Suppose $C$ has a basis $b_1,\ldots,b_n$ and a diagonal idempotent
$\eps$ as in Lemma~\ref{lem:diagonal}. Write
$\eps=\sum_i b_i\otimes c_i$. Then
\begin{equation}\label{eq:trace-duality}
 a=\sum_j\Tr(ab_j)c_j\qquad(a\in C).
\end{equation}
In particular the trace pairing is perfect.
\end{proposition}
\begin{proof}
Let $\lambda_i$ be the basis coordinate maps. Comparing coefficients
in $(z\otimes1)\eps=(1\otimes z)\eps$ gives
\[
 zc_i=\sum_j\lambda_i(zb_j)c_j.
\]
Substitute $z=ab_i$ in the $i$th equation and sum over $i$.
The left side is $a$ because $\sum_i b_ic_i=\mu(\eps)=1$.
The coefficient of $c_j$ on the right is
\[
 \sum_i\lambda_i(ab_i b_j)
   =\sum_i\lambda_i((ab_j)b_i)=\Tr(ab_j),
\]
which proves \eqref{eq:trace-duality}. Applying coordinate maps to
this identity gives an inverse for the square matrix
$(\Tr(b_i b_j))$.
\end{proof}

\begin{lemma}\label{lem:corner-trace}
If $C$ has a finite basis and nondegenerate trace pairing, then every
idempotent factor $eC$ has these properties.
\end{lemma}
\begin{proof}
The image of multiplication by $e$ has a basis obtained by Gaussian
elimination. As $k$-modules,
$C=eC\oplus(1-e)C$. Multiplication by $a\in eC$ acts as zero on
the second summand, so $\Tr_C(a)=\Tr_{eC}(a)$.
If $a\in eC$ pairs to zero with every element of $eC$, then for
every $b\in C$,
\[
 \Tr_C(ab)=\Tr_{eC}(a(eb))=0.
\]
Nondegeneracy on $C$ gives $a=0$.
\end{proof}

\begin{lemma}\label{lem:monogenic-separable}
Let $C$ be a reduced finite-basis $k$-algebra with nondegenerate trace.
If $C=k[x]$ and $f$ is the minimal monic polynomial of $x$, then
$f'$ is invertible in $C\simeq k[T]/(f)$.
\end{lemma}
\begin{proof}
First, every $k$-derivation $\delta:C\to C$ is zero. For $a,b\in C$,
write $m_a$ for multiplication by $a$. The identity
\[
 [b\delta,m_a]=m_{b\delta(a)}
\]
and vanishing of the trace of a commutator give
$\Tr(b\delta(a))=0$. Nondegeneracy gives $\delta(a)=0$.

If $f'(x)h=0$, the assignment $x\mapsto h$ defines a derivation
on $k[T]/(f)$: it sends a polynomial $g$ to $g'(x)h$, and the
relation $f$ maps to zero. Hence $h=0$. Finally, regularity from
Lemma~\ref{lem:reduced-regular} supplies
$f'(x)=f'(x)^2b$. The element $1-f'(x)b$ annihilates $f'(x)$ and
must vanish. Thus $f'(x)b=1$.
\end{proof}

\section{Constructing a tensor separator from finitely many scalars}

Fix a basis $b_1,\ldots,b_n$ of $C$ and its diagonal idempotent
$\eps\in D=C\otimes_k C$. We may assume $n>0$, since $n=0$ gives
the zero algebra. We seek $x\in C$ and $w\in D$ such that
\begin{equation}\label{eq:separator}
 (\Del x)w=1-\eps.
\end{equation}

\begin{lemma}\label{lem:primitive-polynomial}
Let $B$ be a finite-basis algebra over $k$. If the coefficients of
$g\in B[Z]$ generate the unit ideal of $B$, multiplication by $g$
on $B[Z]$ is injective.
\end{lemma}
\begin{proof}
Suppose $gq=0$ for a nonzero $q$. Among nonzero annihilators of degree
at most $\deg q$, choose one of least degree, say $q_0$ of degree $s$.
This is a finite linear-algebra construction: for each degree bound,
the annihilating equation is a finite homogeneous system over $k$.

Write $g=\sum_{i=0}^d a_iZ^i$. Comparing highest coefficients shows
that $a_dq_0$ has degree less than $s$. It is again an annihilator
of $g$, so minimality gives $a_dq_0=0$.
Now $\sum_{i<d}a_iZ^i$ also annihilates $q_0$. Descending through its
coefficients gives $a_iq_0=0$ for every $i$, by the same argument.
A linear combination of the $a_i$ is $1$, so $q_0=0$, a contradiction.
\end{proof}

\begin{proposition}\label{prop:separator}
If $k$ contains $n^3+1$ specified distinct scalars, then
\eqref{eq:separator} has a solution.
\end{proposition}
\begin{proof}
The ideal $\ker\mu$ is generated by the $\Del b_i$, since the $b_i$
span $C$ and therefore generate it as an algebra. Since
$1-\eps\in\ker\mu$, the elements
$\eps,\Del b_1,\ldots,\Del b_n$ generate the unit ideal of $D$.
Apply Lemma~\ref{lem:primitive-polynomial} to
\[
 g(Z)=\eps+\sum_{i=1}^n\Del b_i Z^i\in D[Z].
\]
The algebra $D$ has the basis $b_i\otimes b_j$, of size $n^2$.
Thus multiplication by $g$ is an injective square matrix over $k[Z]$.
Its determinant $h(Z)$ is nonzero: otherwise Gaussian elimination
over $k(Z)$, followed by clearing denominators, would give a nonzero
polynomial vector in its kernel. Each matrix entry has degree at most
$n$, so $\deg h\leq n^3$.

A nonzero polynomial of degree at most $n^3$ cannot vanish at all
the specified $n^3+1$ distinct scalars. Choose $a$ with $h(a)\ne0$.
The adjugate matrix makes multiplication by $g(a)$ invertible, hence
$g(a)$ is a unit of $D$. Put
\[
 x=\sum_{i=1}^n a^i b_i,\qquad
 w=(1-\eps)g(a)^{-1}.
\]
Since $g(a)=\eps+\Del x$, $\eps\Del x=0$ and
$\eps(1-\eps)=0$, these elements satisfy \eqref{eq:separator}.
\end{proof}

\section{A tensor separator gives a global primitive element}

\begin{lemma}\label{lem:reflect-units}
An injective ring map from a regular ring reflects units.
\end{lemma}
\begin{proof}
If $a=a^2b$, then $1-ab$ annihilates $a$. When the image of $a$
is a unit, the image of $1-ab$ is zero. Injectivity gives $ab=1$.
\end{proof}

\begin{proposition}\label{prop:global-generator}
Let $C$ be reduced with a finite $k$-basis and a diagonal idempotent.
If $x,w$ satisfy \eqref{eq:separator}, then $C=k[x]$, and the minimal
monic polynomial of $x$ is separable.
\end{proposition}
\begin{proof}
Let $f$ be the minimal monic annihilator of $x$, of degree $m$.
If $m=0$, its equation is $1=0$, and the assertion is immediate.
Assume $m>0$. There is a polynomial $H(X,Y)$ with
\[
 f(X)-f(Y)=(X-Y)H(X,Y),\qquad H(Y,Y)=f'(Y).
\]
It is monic of degree $m-1$ in $X$. These identities follow from
Horner recursion: the divided difference of $a+TF(T)$ is
$F(Y)+XH_F(X,Y)$.

Evaluate at $X=x\otimes1$, $Y=1\otimes x$, and denote the result
by $H_x$. The root equation gives $(\Del x)H_x=0$. Multiplication
by $w$ yields $(1-\eps)H_x=0$. Equalisation on the diagonal then gives
\begin{equation}\label{eq:divided-diagonal}
 H_x=\eps H_x=\eps(1\otimes f'(x)).
\end{equation}

Minimality makes $1,x,\ldots,x^{m-1}$ linearly independent over $k$.
Gaussian elimination constructs a linear functional $\lambda:C\to k$
with
\[
 \lambda(1)=\cdots=\lambda(x^{m-2})=0,
 \qquad \lambda(x^{m-1})=1.
\]
For $m=1$, the first list of conditions is empty. One explicit
construction is to project onto the quotient by the span of the lower
powers, take a nonzero coordinate of the image of $x^{m-1}$, and
normalise it. Contract the left tensor factor of
\eqref{eq:divided-diagonal} by $\lambda$. Since $H$ is monic in $X$,
this gives
\begin{equation}\label{eq:derivative-inverse}
 1=f'(x)(\lambda\otimes1)(\eps).
\end{equation}
Thus the derivative already has an explicit inverse in $C$.

Let $Q=k[T]/(f)$, viewed as the subalgebra $k[x]\subseteq C$
by Lemma~\ref{lem:minimal}. It has its monomial basis and is reduced,
so it is regular by Lemma~\ref{lem:reduced-regular}.
Lemma~\ref{lem:reflect-units} therefore puts $f'(x)^{-1}$ in $Q$.
In particular $f$ is separable over $k$.

It remains to show that $Q=C$. Write $\eps=\sum_j b_j\otimes c_j$
in an arbitrary basis of $C$. Every right-hand coordinate of $H_x$
in this left basis belongs to $Q$: expand the powers of $x$ in the
left tensor factor in the basis $b_j$; the remaining coefficients
are polynomials in the right-hand $x$. By \eqref{eq:divided-diagonal},
\[
 \eps=H_x(1\otimes f'(x)^{-1}),
\]
so every $c_j$ belongs to $Q$ as well. Finally
\eqref{eq:trace-duality} expresses each $a\in C$ as
$\sum_j\Tr(ab_j)c_j$, an element of $Q$.
\end{proof}

\section{Bounded coefficient closure and the small-field branch}

\begin{lemma}\label{lem:closure}
Given finitely many scalars in $k$ and a natural number $M$, one can
construct either $M+1$ distinct elements of $k$, or an explicitly
finite subfield $k_0\subseteq k$ containing all the given scalars.
\end{lemma}
\begin{proof}
Maintain a finite list of distinct scalars. Test whether it contains
$0,1$ and the given scalars and is closed under addition, negation,
multiplication and total inverse, with $0^{-1}=0$.
Each test is finite and decidable. If it fails, the failure exhibits
one missing scalar, which is appended. If it succeeds, the listed
elements form the required finite subfield. Stop as soon as the list
has $M+1$ elements. Starting with the empty list, at most $M+1$
unsuccessful tests are possible before this alternative is reached.
\end{proof}

\begin{lemma}\label{lem:finite-regular-splitting}
An explicitly finite, reduced algebra over a discrete field is a
finite product of explicitly finite fields, with zero factors allowed.
\end{lemma}
\begin{proof}
The algebra $A$ is regular by Lemma~\ref{lem:reduced-regular}.
For each element $a$ in a finite enumeration choose $b$ with
$a=a^2b$. Set $f=ab$. Then
\[
 f^2=f,\qquad a(1-f)=0.
\]
On the factor $fA$, the element $fa$ has inverse $fb$; on
$(1-f)A$, the image of $a$ is zero.

Process the enumeration successively. A current factor $eA$ is
replaced by $efA$ and $e(1-f)A$. An element already known to be zero
or invertible retains that property on both new factors.
After the finite list has been processed, every element of each
factor is zero or a unit. Equality is decidable, so each factor can
be tested for being the zero ring. Every other factor is a field.
Its finite enumeration is obtained by projecting the original
enumeration and deleting duplicates.

At each stage the maps
\[
 eA\longrightarrow efA\times e(1-f)A,
 \quad z\longmapsto(fz,(1-f)z),
 \qquad (u,v)\longmapsto u+v
\]
are inverse algebra maps. Composing them proves the product assertion.
Termination is induction on the enumeration length, without any
search for a prime ideal.
\end{proof}

\begin{lemma}\label{lem:finite-field-primitive}
An explicitly finite field $F$ is generated by one element over any
subfield $k_0\subseteq F$.
\end{lemma}
\begin{proof}
The multiplicative order of each nonzero element can be found by
finite search in its powers. Choose $\gamma$ of maximal order $m$.
For two commuting elements of orders $r,s$, there is an element of
order $\operatorname{lcm}(r,s)$: factor the positive integers $r,s$,
take for each prime a power of the element with the larger prime-power
order, and multiply these components. Components of coprime orders
have product of order the product, by B\'ezout's identity.
This is a finite calculation with integers and group operations.

Maximality of $m$ now implies that the order of every nonzero element
of $F$ divides $m$. All nonzero elements are therefore roots of
$T^m-1$. The $m$ distinct powers of $\gamma$ already give $m$ roots;
any further element would contradict the root bound for a polynomial
over a field. Thus every nonzero element is a power of $\gamma$,
and every element of $F$ is a polynomial in $\gamma$ over $k_0$.
\end{proof}

\begin{proposition}\label{prop:small}
Suppose $C$ is reduced with a finite basis $b_1,\ldots,b_n$ and
nondegenerate trace. Suppose also that a finite subfield $k_0\subseteq k$
contains the coefficients of monic annihilators of the $b_i$.
Then $C$ is a finite product of separable monogenic $k$-algebras.
\end{proposition}
\begin{proof}
Let $A=k_0[b_1,\ldots,b_n]\subseteq C$. The monic equations provide
a finite spanning list of bounded monomials over $k_0$.
Since $k_0$ is explicitly finite, the corresponding finite list of
linear combinations enumerates $A$. Equality in $C$, and hence in
$A$, is decidable from the given $k$-basis. The subalgebra $A$ is
reduced, so Lemma~\ref{lem:finite-regular-splitting} supplies its
finite idempotent decomposition.

For a nonzero piece $eA$, choose a primitive element $\gamma$ by
Lemma~\ref{lem:finite-field-primitive}. The inclusion $eA\subseteq eC$
preserves the identity $e$. Each $eb_i$ is a polynomial in $\gamma$
over $k_0$, with coefficients interpreted as multiples of $e$.
Since the $eb_i$ span $eC$ over $k$, this gives
\[
 eC=k[\gamma].
\]
Lemma~\ref{lem:corner-trace} supplies a finite basis and nondegenerate
trace for $eC$, and reducedness is inherited from $C$.
Its minimal polynomial is separable by
Lemma~\ref{lem:monogenic-separable}. Thus $eC\simeq k[T]/(f_e)$
with $f_e$ monic separable. If $eA=0$, injectivity into $C$ gives
$e=0$ and $eC=0$, represented by $f_e=1$.

The idempotents sum to $1$ in $A$, hence also in $C$. The same
projection and addition maps assemble
$C\simeq\prod_e eC$, with the required monogenic presentations.
\end{proof}

\section{Completion of the theorem}

\begin{proof}[Proof of Theorem~\ref{thm:main}]
Construct $\eps$ by Lemma~\ref{lem:diagonal}. Proposition~\ref{prop:regular}
makes $C$ regular and reduced. Proposition~\ref{prop:finite} constructs
a finite basis $b_1,\ldots,b_n$, and Proposition~\ref{prop:trace}
gives a perfect trace pairing. If $n=0$, take the quotient by $1$.

Otherwise choose monic annihilators of the $b_i$ and collect their
finitely many coefficients. Apply Lemma~\ref{lem:closure} with
$M=n^3$. In the first alternative, Proposition~\ref{prop:separator}
gives a tensor separator, and Proposition~\ref{prop:global-generator}
gives a single separable monogenic presentation of $C$.
In the second alternative, Proposition~\ref{prop:small} gives the
finite product. This proves the asserted decomposition.

For completeness, a monic separable quotient $k[T]/(f)$ is finite
free and formally \'etale. To lift a root across a square-zero ideal
$N$ of an algebra $B$, choose a lift $t$ of the root. The derivative
$f'(t)$ is a unit, because it is a unit modulo $N$. The unique
correction in $N$ is
\[
 t\longmapsto t-f'(t)^{-1}f(t),
\]
as $f(t+u)=f(t)+f'(t)u$ for $u\in N$. A unit modulo a square-zero
ideal lifts to a unit by inverting $1+v$ as $1-v$ for $v\in N$.
Finite products preserve this lifting property: their component
idempotents lift uniquely across a square-zero ideal, after which
one lifts each component map. Together with finite presentation,
this proves that the product is finite \'etale.
\end{proof}

\begin{remark}
The small-field alternative is essential. In $\mathbb F_2^3$, every
element satisfies $T^2-T$, so the subalgebra generated by one element
has dimension at most two. The product decomposition nevertheless
has three monogenic factors. Over an arbitrary discrete field the
argument supplies separable polynomials without requiring their
factorisation into irreducibles. If such factorisations are supplied,
the Chinese remainder theorem refines the result to a product of
finite separable field extensions.
\end{remark}

\part{Connection with the formalisation}

\section{The precise result in Unit}

The mathematical product theorem above is formalised in
\path{unit/sag/etale-field-product}. Its endpoint is
\begin{center}
\texttt{UnramifiedFieldProduct.Over[k,setk,field]}\\
\texttt{.At[C,fp,unique].theorem}.
\end{center}
Here \texttt{field} supplies the zero-or-unit operation and $1\ne0$;
\texttt{fp} is a finite presentation under propositional truncation;
and \texttt{unique} is precisely uniqueness of square-zero lifting.
Neither a basis nor a specified prime is an input.

The conclusion is the propositional truncation of
\texttt{ProductPresentation.Over[k,setk].At[C]}.
This record contains a finite binary tree of monic polynomials,
an inverse of each derivative in the corresponding quotient, and
an algebra isomorphism from the actual product to $C$.
Both directions are algebra maps with proved inverse equations.
The polynomials need not be irreducible, and zero factors are retained.
The product is represented by its finite tree rather than flattened
to one list of factors.

Existence statements in the mathematical argument are expressed
through propositional truncation where appropriate. Finite witnesses
are combined inside that truncation. There is no global choice of
representatives or primitive elements. The development uses the
existing univalence, replacement and surjection parameters of the
shared theory and introduces no new axiom.

\section{Correspondence of the proof steps}

\begin{longtable}{@{}>{\raggedright\arraybackslash}p{0.30\textwidth}
 >{\raggedright\arraybackslash}p{0.66\textwidth}@{}}
\textbf{Mathematical step} & \textbf{Principal formal components}\\
\hline
\endfirsthead
\textbf{Mathematical step} & \textbf{Principal formal components}\\
\hline
\endhead
Diagonal and flatness transfer &
\texttt{DiagonalCertificate}, \texttt{FlatDiagonalTransfer},
\texttt{UnramifiedFieldRegular}. These construct the diagonal and
the averaging retraction from formal unramifiedness.\\[5pt]
Finite span and basis &
\texttt{ZeroDimensionalFinite},
\texttt{UnramifiedFieldFiniteGeneral},
\texttt{FinitePresentedFieldBasis}. Monic substitutions, integral
descent, the finite monic cover and its relation matrix are constructed.\\[5pt]
Trace duality &
\texttt{TensorBasisConstruct}, \texttt{DiagonalTraceBridge},
\texttt{TraceDuality}, \texttt{FiniteUnramifiedTrace}. The tensor
coordinates and their reconstruction are proved from the coproduct.\\[5pt]
Minimal monic quotient &
\texttt{MinimalPolynomial}, \texttt{MinimalPolynomialInjection}.
The Gaussian search, injection and inverse algebra map when the
generator is surjective are explicit.\\[5pt]
Monogenic separability &
\texttt{TraceDerivations}, \texttt{TraceRetraction},
\texttt{MonicTraceSeparable}, \texttt{GeneratedSeparableMonogenic}.
The derivation killing $f'$ is constructed using dual numbers.\\[5pt]
Tensor separator &
\texttt{DiagonalSeparator}, \texttt{FinitePolynomialSelection},
and the polynomial-basis and determinant modules. The scalar search
has the explicit bound $n^3+1$.\\[5pt]
Leading-coefficient functional &
\texttt{FieldSeparatingFunctional}, \texttt{MinimalPowerFunctional}.
These construct the quotient coordinate functional used in
\eqref{eq:derivative-inverse}.\\[5pt]
Divided difference and derivative inverse &
\texttt{MonicDividedDifference}, \texttt{TensorDividedDifference},
\texttt{TensorMonicContraction}, \texttt{SeparatorDerivativeUnit}.\\[5pt]
Global generation &
\texttt{RegularUnitReflection}, \texttt{TensorPolynomialCoordinates},
\texttt{SeparatorMonogenic}, \texttt{LargeFieldMonogenic}.
Preimages of the diagonal coordinates and then of every element
are constructed.\\[5pt]
Finite coefficient alternative &
\texttt{FiniteMonicCoefficientFamily}, \texttt{FiniteCoefficientField},
\texttt{MonicFiniteImage}. The finite coefficient field is constructed,
and the monic equations are lifted to it.\\[5pt]
All finite idempotent pieces &
\texttt{FiniteRegularRefinement}, \texttt{FiniteResolvedField}.
Both branches of every split are retained; induction is on the
enumeration length.\\[5pt]
Primitive elements on the pieces &
\texttt{FiniteFieldPrimitive}, \texttt{HornerCoefficientChange},
\texttt{PolynomialSpanGeneration}, \texttt{IdempotentMonogenicSpan}.
Polynomial representatives span each ambient corner.\\[5pt]
Product assembly &
\texttt{IdempotentProductIsomorphism}, \texttt{ProductsOnIdempotents},
\texttt{SmallFieldProduct}, \texttt{FiniteEtaleProduct},
\texttt{UnramifiedFieldProduct}.\\
\end{longtable}

There are two useful distinctions between the mathematical presentation
and its implementation. First, the determinant argument in
Proposition~\ref{prop:separator} is implemented using matrix McCoy
rather than constructing $k(Z)$ and clearing denominators. Polynomial
McCoy supplies the preceding regularity statement. Second, trace
nondegeneracy is passed to an idempotent factor using rectangular
trace cyclicity and a retraction, without constructing a basis of the
complementary factor. Both implementations establish exactly the
lemmas used above.

The global-generator proof is the shorter argument used in the new
product development. The earlier selected-prime development instead
used relative bases on an idempotent neighbourhood of a given prime.
Those earlier files remain available, but the product theorem has no
prime parameter and does not use that selection step.

\section{Validation and the remaining local connection}

The recorded fresh check of the complete import tree gives
$18757$ declarations, $0$ failures, in $9.42$ seconds.
Of these, $847$ declarations in $25$ owned modules are new.
All $446$ modules pass the module checks. The opening audit covers
all $797$ named new declarations and reports no openings outside
their declared lists; each new declaration takes at most $0.03$
seconds and opens at most $18$ definitions in that run.

Regressions cover $\mathbb F_2$ from finite presentation and formal
unramifiedness, the constructed standard algebra for $T^2-T$,
the large branch with its actual tensor separator, the small branch
on $\mathbb F_2^3$, and the zero algebra. The $\mathbb F_2^m$
regression constructs its basis, trace pairing and monic equations
and invokes the full small-field construction. It does not supply the
desired product isomorphism. These checks validate proof terms;
they do not extract and normalise the entire returned truncated tree.
Three isolated mutations---an incorrect divided difference, an
incorrect complementary piece, and a missing derivative inverse---were
rejected by the checker.

The formal endpoint is the explicit product isomorphism with derivative
inverses. The final paragraph of the mathematical proof explains why
these algebras are finite \'etale; the product milestone does not add
a separate packaged theorem identifying its conclusion with every
existing \'etale interface.

For the local structure problem over a general base, this result
provides the full residue-field decomposition. Connecting it to the
original local factor-selection proof still requires assembling the
coefficient and denominator bridge on the exact special fibre and
supplying the remaining local hypotheses. The existing selected-factor
local adapter is unchanged. In particular, completion
of this field theorem does not by itself assert completion of the
unconditional local theorem that flat and unramified implies \'etale.

The implementation signatures, commands and validation details are
recorded in \path{unit/sag/etale-field-product/README.txt}.
The earlier note \path{etale-field-reduction.tex} also describes the
selected-factor route and the local integration work; the present
version separates the mathematical product proof from these formal
connections.

\end{document}
